\documentclass[11pt,a4paper]{article}
\usepackage[T1]{fontenc}
\usepackage[utf8]{inputenc}
\usepackage{lmodern,amsmath,amssymb}
\usepackage[margin=25mm]{geometry}
\usepackage[hidelinks]{hyperref}
\hypersetup{pdftitle={The Principia's fifth postulate: joint determinacy of place and motion},pdfauthor={navstokgap research working notes}}
\newcommand{\tightlist}{\setlength{\itemsep}{0pt}\setlength{\parskip}{0pt}}
\setlength{\emergencystretch}{3em}
\setcounter{secnumdepth}{0}
\begin{document}
\hypertarget{the-principias-fifth-postulate-joint-determinacy-of-place-and-motion}{%
\section{The Principia's fifth postulate: joint determinacy of place and
motion}\label{the-principias-fifth-postulate-joint-determinacy-of-place-and-motion}}

\textbf{Result, 2026-09-27 (user's goal: identify the statement whose
change gives mechanics with \(h>0\), as the parallel postulate does for
curved geometry).} The statement is \textbf{joint determinacy}: the
place of a body and its quantity of motion are quantities given
together, so that the observables of a body multiply commutatively.
Newton never states it as a postulate. He uses it in Definition II, in
Lemma X and in the polygon of Proposition I, and its explicit trace is
the last sentence of the scholium closing Book I, Section I: the
quantities compared are ``always to be diminished without limit''. Four
theorems make the analogy with Euclid's fifth postulate exact.

\begin{enumerate}
\def\labelenumi{\arabic{enumi}.}
\tightlist
\item
  \textbf{Independence (Theorem A).} For every value of a real constant
  \(\hbar\), the Moyal product \(*_\hbar\) on the observables of one
  degree of freedom satisfies Newton's second law as an identity of
  observables, and the flows of Newton's own integrable cases (inertia,
  uniform gravity, force proportional to distance) are the classical
  flows and automorphisms of the product. The Laws cannot decide joint
  determinacy, which holds exactly when \(\hbar=0\).
\item
  \textbf{Exactly one constant (Theorem B).} Every associative product
  on the observables that deforms the commutative one, is covariant
  under the group generated by those three flows (the affine symplectic
  group), and respects complex conjugation, is a Moyal product
  \(*_\hbar\) with one real constant \(\hbar\), which has units of
  action.
\item
  \textbf{The floor (Theorem C).} For \(\hbar\ne0\), every regular state
  obeys \(\Delta q\,\Delta p\ge|\hbar|/2\), and every recorded
  comparison of the inertial line with the constant-force parabola obeys
  the floor of the \href{planck-gap-paper.md}{Planck paper}, Theorem 6:
  \(\frac s8\sum_j\Delta(\hat D_j)+\frac J2\sum_j\Delta(\hat X_j)\ge|\hbar| \arcsin(1-2\epsilon)\).
  For \(\hbar=0\) there is no floor (Theorem I of the
  \href{necessity-unit-and-indeterminacy.md}{unit-and-indeterminacy
  note}).
\item
  \textbf{The mechanical Wallis postulate (Theorem D).} The similarity
  that rescales every action, \((q,p)\mapsto(\lambda q,\lambda p)\), is
  a symmetry of \(*_\hbar\) exactly when \(\hbar=0\). Wallis showed that
  Euclid's fifth postulate is equivalent to the existence of similar
  figures of different size; here joint determinacy is equivalent to
  mechanical similarity of recorded motions at every scale, and
  \(\hbar\) plays the part of the absolute length of non-Euclidean
  geometry. Up to isomorphism there are exactly two structures,
  \(\hbar=0\) and \(\hbar\ne0\).
\end{enumerate}

The mathematics is established: Moyal's product
(\href{https://doi.org/10.1017/S0305004100000487}{Moyal 1949},
metadata), Groenewold's observation that its bracket reduces to the
Poisson bracket for quadratic functions
(\href{https://doi.org/10.1007/978-94-017-6065-2_1}{Groenewold 1946},
metadata), the deformation-quantization programme
(\href{https://doi.org/10.1016/0003-4916(78)90224-5}{Bayen, Flato,
Fronsdal, Lichnerowicz and Sternheimer 1978}, metadata), and Robertson's
inequality (\href{https://doi.org/10.1103/PhysRev.34.163}{Robertson
1929}, metadata). The uniqueness in Theorem B is Gutt's theorem: the
Moyal star product is the unique
\(Sp(2n,\mathbb R)\ltimes\mathbb R^{2n}\)-invariant and covariant star
product on \(\mathbb R^{2n}\) (S. Gutt, \emph{Mém. Acad. Roy. Belg. Cl.
Sci.} 44:6, 1983, as cited in
\href{https://doi.org/10.1007/s00220-003-0973-7}{Duval, El Gradechi and
Ovsienko 2004}, passage). The proof below is a short self-contained
version for the products considered here. The contribution is the
identification: which statement of the \emph{Principia} is the one to
change, and the proof that the Laws are independent of it in Newton's
own integrable cases.

\hypertarget{the-textual-anchor}{%
\subsection{1. The textual anchor}\label{the-textual-anchor}}

The scholium closing Book I, Section I (1687 text, held in this
repository's
\href{../docs/classics/Newton_Principia_1687_BookI_SectI_scholia_la_wikisource.md}{Latin
companion}), ends:

\begin{quote}
Igitur in sequentibus, siquando facili rerum imaginationi consulens,
dixero quantitates quam minimas, vel evanescentes vel ultimas, cave
intelligas quantitates magnitudine determinatas, sed cogita semper
diminuendas sine limite.
\end{quote}

``So in what follows, if for ease of imagination I speak of quantities
as least, or vanishing, or ultimate, do not understand quantities
determinate in magnitude, but think of them always as to be diminished
without limit.'' Newton defends this against the objection from Euclid,
\emph{Elements} X, and applies it to the sagitta of Lemma X and the
chords of Proposition I. The sentence is about geometry; its application
to motion needs one more, unstated, premise: that the place of the body
and its quantity of motion (Definition II) are given together at each
instant, so that the sagitta measured from a place and the impulse
inferred from a change of motion can both be diminished without limit in
one and the same comparison. That premise is joint determinacy. Its
algebraic form is the commutativity of the product of observables.

The other division scholium, after Lemma XI, says of contact angles that
between any two a new series can be inserted, ``Neq; novit natura
limitem''. That statement concerns unobserved insertion, and it remains
true for every \(\hbar\): the
\href{refinement-composition-and-limit.md}{refinement note} shows that
inserting unobserved times composes exactly at every \(\hbar\). The
change is confined to recorded comparisons.

\hypertarget{the-family-of-products}{%
\subsection{2. The family of products}\label{the-family-of-products}}

Observables of one degree of freedom (the transverse coordinate of
Galileo's comparison) are polynomials in \(q\) and \(p\) with complex
coefficients. Write \(\mu(f\otimes g)=fg\) and

\[P=\partial_q\otimes\partial_p-\partial_p\otimes\partial_q,\qquad
f*_\hbar g=\mu\circ\exp\Bigl(\frac{i\hbar}2P\Bigr)(f\otimes g),\]

a finite sum on polynomials. For \(\hbar=0\) it is the commutative
product; \(q*_\hbar p-p*_\hbar q=i\hbar\). The Moyal bracket is
\(\{f,g\}_\hbar=(f*_\hbar g-g*_\hbar f)/(i\hbar)=\frac2\hbar\, \mu\circ\sin(\frac\hbar2P)(f\otimes g)\),
with \(\{f,g\}_0\) the Poisson bracket.

\hypertarget{theorem-a-independence}{%
\subsection{3. Theorem A: independence}\label{theorem-a-independence}}

\textbf{Theorem A.} Let \(H=p^2/(2M)+V(q)\) with \(V\) a polynomial and
\(M>0\). For every real \(\hbar\):

\begin{enumerate}
\def\labelenumi{(\alph{enumi})}
\item
  \(\{q,H\}_\hbar=p/M\) and \(\{p,H\}_\hbar=-V'(q)\): the Heisenberg
  equations are Newton's second law, \(M\ddot q=-V'(q)\), as an identity
  of observables;
\item
  if \(\deg V\le2\), which covers inertia (\(V=0\)), uniform gravity
  (\(V=-Fq\)) and the force proportional to distance of Book I,
  Proposition X (\(V=kq^2/2\)), then \(\{f,H\}_\hbar=\{f,H\}_0\) for
  every \(f\): every observable evolves by the classical flow
  \(\varphi_t\), and each \(\varphi_t\) is an automorphism of
  \(*_\hbar\).
\end{enumerate}

So the commutative product (\(\hbar=0\)) and every Moyal product
(\(\hbar\ne0\)) satisfy the Laws and reproduce Newton's integrable
motions exactly; joint determinacy holds in the first and fails in the
others.

\emph{Proof.} The bracket \(\{f,g\}_\hbar\) is the Poisson bracket plus
terms \(\hbar^{2j}\,\mu\,P^{2j+1}(f\otimes g)\), \(j\ge1\), each of
which puts at least three derivatives on each factor. If one factor has
degree at most two they vanish. This gives (a) because \(q\) and \(p\)
have degree one, and (b) because \(H\) has degree two. For (b), the flow
of a quadratic Hamiltonian is an affine symplectic map, and \(*_\hbar\)
is covariant under affine symplectic maps: translations commute with the
constant-coefficient operator \(P\), and a linear map of determinant one
leaves \(P\) invariant. \(\square\)

\hypertarget{theorem-b-exactly-one-action-constant}{%
\subsection{4. Theorem B: exactly one action
constant}\label{theorem-b-exactly-one-action-constant}}

Consider products on \(\mathbb C[q,p]\) of the form
\(f*g=\mu\circ B(f\otimes g)\), with
\(B=\sum_{\alpha,\beta}c_{\alpha\beta}(q,p)\, \partial^\alpha\otimes\partial^\beta\)
a bidifferential operator (a finite sum on each pair of polynomials),
such that

\begin{itemize}
\tightlist
\item
  (H1) \(*\) is associative with unit 1, and \(c_{00}=1\) (it deforms
  the commutative product);
\item
  (H2) \(*\) is covariant under the affine symplectic group: for every
  translation and every linear map \(A\) of determinant one,
  \((f\circ A)*(g\circ A)=(f*g)\circ A\);
\item
  (H3) \(\overline{f*g}=\bar g*\bar f\).
\end{itemize}

By Theorem A(b), the flows of inertia, uniform gravity and the force
proportional to distance are affine symplectic maps; free flight is the
shear \((q,p)\mapsto(q+tp/M,p)\), uniform gravity adds translations, and
the harmonic flow is a rotation in suitable units. The shears and
rotations generate the group of linear maps of determinant one. So (H2)
says that the motions Newton integrated in closed form are symmetries of
the product.

\textbf{Theorem B.} Under (H1)--(H3), \(*=*_\hbar\) for a unique real
\(\hbar\). The constant \(\hbar\) has the units of action, the
reciprocal of the units of \(P\).

\emph{Proof.} Translation covariance makes the coefficients
\(c_{\alpha\beta}\) constant. Then \(B\) is a polynomial (a formal
series, acting finitely on polynomials) in the components of two vectors
\(u=\partial^{(1)}\) and \(v=\partial^{(2)}\), invariant under the
simultaneous action of the group of determinant-one linear maps. By the
first fundamental theorem of invariant theory for \(SL(2)\), such
invariants are functions of the single determinant \(\det(u,v)=P\). So
\(B=F(P)\) with \(F\) a formal power series, \(F(0)=1\). For three
factors write \(x=P_{12}\), \(y=P_{13}\), \(z=P_{23}\) for the
determinants of the pairs of derivative vectors. The three \(2\times2\)
minors of a \(2\times3\) matrix are algebraically independent, and
\((f*g)*h\), \(f*(g*h)\) correspond to \(F(x)F(y+z)\) and
\(F(z)F(x+y)\). Associativity therefore gives \(F(x)F(y+z)=F(z)F(x+y)\)
as formal series; setting \(z=0\) gives \(F(x)F(y)=F(x+y)\), so
\(F(x)=e^{\kappa x}\). Exchanging the two factors reverses the sign of
\(P\) and complex conjugation leaves \(P\) real, so (H3) gives
\(\bar F(x)=F(-x)\), that is \(\bar\kappa=-\kappa\), and
\(\kappa=i\hbar/2\) with \(\hbar\) real. Units: \(P\) carries the
reciprocal of the units of \(q\,p\). \(\square\)

\textbf{Only two structures up to isomorphism.} The dilation
\(D_\lambda(q,p)=(\lambda q,\lambda p)\) satisfies
\(P((f\circ D_\lambda)\otimes(g\circ D_\lambda))=\lambda^2(P(f\otimes g))\circ D_\lambda\),
so \(f\mapsto f\circ D_\lambda\) is an isomorphism from
\(*_{\lambda^2\hbar}\) to \(*_\hbar\), and \((q,p)\mapsto(q,-p)\)
carries \(*_\hbar\) to \(*_{-\hbar}\) (with complex conjugation). All
products with \(\hbar\ne0\) are therefore isomorphic, and distinct from
the commutative one. Geometry has three constant-curvature classes
(flat, spherical, hyperbolic); mechanics has two. The value of \(\hbar\)
in physical units is fixed by an external standard, as the absolute
length of non-Euclidean geometry is.

\hypertarget{theorem-c-the-floor-on-the-recorded-comparison}{%
\subsection{5. Theorem C: the floor on the recorded
comparison}\label{theorem-c-the-floor-on-the-recorded-comparison}}

A state is a linear functional \(\omega\) on the observables with
\(\omega(1)=1\) and \(\omega(\bar f*_\hbar f)\ge0\); it is regular if
its GNS representation integrates to Weyl operators.

\textbf{Theorem C.} (a) For \(\hbar\ne0\) every state satisfies
\(\Delta q\,\Delta p\ge|\hbar|/2\). (b) For \(\hbar\ne0\) and regular
states, every recorded comparison of the inertial line with the
constant-force parabola, over a cell of duration \(\tau\) with fall
\(s=F\tau^2/(2M)\) (the Planck paper's convention) and impulse
\(J=F\tau\), made by any instrument at error probability
\(\epsilon<\frac12\), leaves undetermined impulses \(\hat D_j\) and
displacements \(\hat X_j\) with

\[\frac s8\sum_j\Delta(\hat D_j)+\frac J2\sum_j\Delta(\hat X_j)\ \ge\
|\hbar|\arcsin(1-2\epsilon)>0 .\]

\begin{enumerate}
\def\labelenumi{(\alph{enumi})}
\setcounter{enumi}{2}
\tightlist
\item
  For \(\hbar=0\), with the instruments of Theorem I of the
  unit-and-indeterminacy note available, the body's posterior area can
  be made arbitrarily small: there is no floor.
\end{enumerate}

\emph{Proof.} (a) The form \((f,g)\mapsto\omega(\bar f*_\hbar g)\) is
positive semidefinite; apply Cauchy--Schwarz to \(f=q-\omega(q)\),
\(g=p-\omega(p)\) and use \(q*_\hbar p-p*_\hbar q=i\hbar\) (Robertson's
argument). (b) A regular state's GNS representation satisfies the Weyl
relations with constant \(\hbar\); by the Stone--von Neumann theorem it
is a multiple of the Schrödinger representation. Theorem 6 of the Planck
paper is proved there from the Weyl relations and Mandelstam--Tamm in
the Bures angle, with \(\hbar\) as the constant, and applies with
\(|\hbar|\) after the reflection \((q,p)\mapsto(q,-p)\) if \(\hbar<0\).
(c) is Theorem I. \(\square\)

Newton's scholium asks that the compared quantities be ``always
diminished without limit''. Theorem A says his Laws allow it or forbid
it equally; Theorem C says that in every model other than the
commutative one, the recorded comparison of the inertial line with the
parabola cannot be diminished below the single constant of Theorem B.

\hypertarget{theorem-d-the-mechanical-wallis-postulate}{%
\subsection{6. Theorem D: the mechanical Wallis
postulate}\label{theorem-d-the-mechanical-wallis-postulate}}

\textbf{Theorem D.} For \(\lambda\ne\pm1\), the action-rescaling
similarity \(D_\lambda\) is an automorphism of \(*_\hbar\) if and only
if \(\hbar=0\).

\emph{Proof.} By §4, \(D_\lambda\) carries \(*_{\lambda^2\hbar}\) to
\(*_\hbar\), and \(*_{\lambda^2\hbar}=*_\hbar\) iff
\(\lambda^2\hbar=\hbar\). \(\square\)

Wallis replaced Euclid's fifth postulate by the existence of a triangle
similar to a given one at any size; the non-Euclidean geometries are
exactly those without similar figures of different size, and they
possess an absolute unit of length. Theorem D is the mechanical
counterpart. Mechanical similarity that rescales actions, the principle
behind comparing motions at every scale, holds exactly when joint
determinacy does; its failure is the existence of an absolute action.
The \href{action-unit-dimensional-selection.md}{dimensional note} states
the same fact from the other side: an admitted similarity that rescales
the observable forces the floor to be zero (its Theorem B).

\hypertarget{what-remains-and-consequence-for-state}{%
\subsection{7. What remains, and consequence for
STATE}\label{what-remains-and-consequence-for-state}}

The theorems are exact; the identification of joint determinacy as the
statement to change is interpretive, because Newton used the premise
without stating it. Two limits of scope. Theorem B concerns one degree
of freedom and bidifferential products; several degrees of freedom
follow with \(Sp(2n)\) in place of \(SL(2)\), whose invariants of two
vectors are again generated by the symplectic pairing. And the
independence of Theorem A is independence from the Laws and from
Newton's integrable cases; it leaves open which physical premise, stated
before quantum mechanics, forces \(\hbar\ne0\). That is STATE item 2,
where Theorem U supplies the unit's value from radiation thermodynamics
and Theorem I says what must be denied. This note shows that the denial
of joint determinacy is, up to isomorphism, a single step with a single
constant.
\end{document}
